\chapter{Contexto de la CMF R2ET: Grupo~1}
\label{ch:contexto}

\section{Descripción del sistema}

La celda del Grupo~1 contiene una máquina de proceso, un robot de transferencia y dos esteras transportadoras de salida E1 y~E2. El robot retira las piezas terminadas de la máquina y las deposita, de forma alternada, en una de las dos esteras. Cada estera tiene longitud para alojar un máximo de tres piezas simultáneas; si se supera ese límite, la producción se bloquea \cite{groover2015automation,benitez2022automatizacion}.

\begin{figure}[H]
\centering
\begin{tikzpicture}[>=Latex, every node/.style={font=\small}]
  \tikzset{
    mach/.style={draw, fill=UniLight, rounded corners=2mm, minimum width=22mm, minimum height=14mm, align=center},
    robot/.style={draw, fill=UniPrimary!15, rounded corners=3mm, minimum width=18mm, minimum height=18mm, align=center, draw=UniPrimaryDeep, line width=1pt},
    belt/.style={draw, fill=UniLight!60, rounded corners=1mm, minimum width=40mm, minimum height=12mm, align=center},
    arr/.style={->, line width=1.0pt, color=UniPrimaryDeep},
    lbl/.style={font=\scriptsize, fill=white, inner sep=1pt},
  }
  \node[mach]  (M)  at (0, 0)    {Máquina\\proceso};
  \node[robot] (R)  at (3.5, 0)  {Robot\\R};
  \node[belt]  (E1) at (7.5, 1.2){Estera E1\\(cap.~3 piezas)};
  \node[belt]  (E2) at (7.5,-1.2){Estera E2\\(cap.~3 piezas)};
  \draw[arr] (M.east) -- node[lbl,above]{pieza} (R.west);
  \draw[arr] (R.east) -- node[lbl,above,sloped]{alt.} (E1.west);
  \draw[arr] (R.east) -- node[lbl,below,sloped]{alt.} (E2.west);
  \draw[arr] (E1.east) -- ++(1.0,0) node[right]{\scriptsize salida};
  \draw[arr] (E2.east) -- ++(1.0,0) node[right]{\scriptsize salida};
\end{tikzpicture}
\caption{Esquema de la CMF R2ET: robot que alimenta dos esteras de salida de forma alternada.}
\label{fig:cmf_r2et}
\end{figure}

\section{Problema de control}

La alternancia del robot por sí sola no resuelve el problema de control. La decisión crítica es cuánto flujo asignar a cada estera y a qué velocidad descargarlas. El problema se complica por tres factores: la saturación, que obliga a redirigir el flujo antes de agotar el margen de capacidad; el vaciado, que corta la línea de salida si la descarga supera de forma sostenida al flujo de entrada; y el cambio de peso de las piezas, que reduce la eficiencia de descarga hasta un $35\%$ en E1 y un $17.5\%$ en E2. El controlador debe detectar esa caída y reaccionar (acelerar la estera, recortar la admisión o desviar el robot hacia E2) sin esperar a que el nivel alcance el umbral de alarma.

\section{Modelo dinámico}

\subsection*{Planta continua: función de transferencia (Guía VIU)}

La función de transferencia continua asignada al Grupo~1 (R2ET) en la guía de la actividad es:

\begin{equation}
G_p(s) = \frac{0.174\,s + 0.3744}{s^2 + 0.785\,s + 0.3744}
\label{eq:planta_continua}
\end{equation}

Esta es una planta de segundo orden con un cero real. Sus características principales se obtienen directamente de los coeficientes:

\begin{itemize}
  \item Frecuencia natural no amortiguada: $\omega_n = \sqrt{0.3744} = 0.612$~rad/s
  \item Factor de amortiguamiento: $\zeta = \dfrac{0.785}{2\,\omega_n} = \dfrac{0.785}{1.224} = 0.641$ \quad (subamortiguado; oscilación moderada y bien amortiguada)
  \item Polos complejos conjugados:
        \[s_{1,2} = -\zeta\omega_n \pm j\,\omega_n\sqrt{1-\zeta^2}
                  = -0.393 \pm j\,0.470 \;\text{rad/s}\]
  \item Cero real: $s_z = -0.3744/0.174 = -2.15$~rad/s \quad ($5.5\times$ más rápido que los polos dominantes)
  \item Ganancia estática: $G_p(0) = 0.3744/0.3744 = 1.0$ (ganancia unitaria en DC)
\end{itemize}

\subsection*{Reducción al modelo de primer orden para el lazo servo (Fase~1)}

Para el lazo local de seguimiento (Fase~1) se aproxima $G_p(s)$ por su constante de tiempo dominante, \emph{conservando la ganancia DC unitaria} de la planta asignada. La aproximación es válida porque:
\begin{enumerate}
  \item El cero ($s_z=-2.15$) está $5.5\times$ más alejado del origen que los polos dominantes; su influencia sobre la respuesta lenta es secundaria.
  \item Los polos complejos están bien amortiguados ($\zeta=0.641$): la respuesta es solo levemente oscilatoria y queda bien capturada por la constante de tiempo dominante $\tau_d = 1/|\operatorname{Re}(s_{1,2})| = 1/0.393 \approx 2.55$~s.
\end{enumerate}

El modelo reducido continuo, que preserva la ganancia estática de $G_p(s)$, es:
\begin{equation}
G_1(s) = \frac{K_f}{\tau_d\,s + 1}, \qquad K_f = G_p(0) = 1.0,\quad \tau_d = 2.55~\text{s}
\label{eq:planta_reducida}
\end{equation}

La figura siguiente compara la respuesta al escalón de $G_p(s)$ y de $G_1(s)$. La reducción reproduce la ganancia final y la escala temporal dominante, aunque simplifica el transitorio de la planta completa (error máximo de $19.3\%$ en la respuesta normalizada). Por ello se usa como banco comparativo de servorregulación en Fase~1; la coordinación de la celda se evalúa aparte con el balance de masa R2ET.

\artifactfigure{figures/plots/fig_reduccion_planta.png}{0.82}{Respuesta al escalón de la planta completa asignada al Grupo~1 frente al modelo reducido de primer orden usado en Fase~1.}

La perturbación de carga (cambio de peso, que reduce la eficiencia de descarga) se modela como un término aditivo $-K_d\,d(t)$ sobre la dinámica.

\subsection*{Discretización ZOH: planta servo (Fase~1)}

El período de muestreo de la planta servo se fija por el criterio de doce muestras por constante de tiempo dominante: con $\tau_d \approx 2.55$~s, $T_s = \tau_d/12 \approx 0.2125$~s. Conviene distinguir dos escalas temporales: este $T_s$ es el de la dinámica servo local de la Fase~1, mientras que la coordinación de la celda R2ET (Fase~2) y las capas de Industria~4.0 operan a un ciclo nominal mayor, de $1$~s. La discretización por retenedor de orden cero (ZOH) de un primer orden da el polo discreto $a = e^{-T_s/\tau_d}$, con $b = K_f(1-a)$ y $b_d = K_d(1-a)$:

\begin{equation}
\boxed{y(k+1) = \underbrace{0.92}_a\,y(k) + \underbrace{0.08}_b\,u(k) - \underbrace{0.024}_{b_d}\,d(k)}
\label{eq:planta_servo}
\end{equation}

\begin{center}
\small
\begin{tabular}{@{}lll@{}}
\toprule
Parámetro & Valor & Origen \\ \midrule
$a = e^{-T_s/\tau_d}$ & $0.92$ & polo discreto ($T_s/\tau_d = 1/12$, $e^{-1/12}=0.920$) \\
$b = K_f(1-a)$ & $0.08$ & ganancia de control ($K_f = G_p(0) = 1.0$) \\
$b_d = K_d(1-a)$ & $0.024$ & sensibilidad a perturbación ($K_d = 0.30$) \\
$u_{ff} = r\,(1-a)/b$ & $0.60$ & feedforward para $r=0.6$ (ganancia DC unitaria $\Rightarrow u_{ss}=r$) \\
\bottomrule
\end{tabular}
\end{center}

Como la ganancia DC es unitaria, el régimen estacionario exige $u_{ss}=r$; se regula a un nivel de operación normalizado $r=0.6$ ($60\%$ del rango), que deja margen de actuador en ambos sentidos para el rechazo de la perturbación de carga.

\subsection*{Modelo R2ET completo (Fase~2)}

En la Fase~2 se simula el balance de masa de las dos esteras. Primero se calcula el estado físico no recortado:

\begin{equation}
\tilde n_i(k+1) = n_i(k)
  + \underbrace{\lambda(k)\,a(k)\,r_i}_{\text{entrada}}
  - \underbrace{\eta\,u_i(k)\,\bigl[1-\ell_i(k)\bigr]}_{\text{salida}}
\label{eq:r2et_raw}
\end{equation}

El estado usado por la simulación solo se satura por abajo, porque no puede haber ocupación negativa:
\begin{equation}
n_i(k+1)=\max\!\left(0,\tilde n_i(k+1)\right)
\label{eq:r2et_aware}
\end{equation}

donde $\lambda(k)$ es el caudal de la máquina, $a(k)$ la fracción de admisión, $r_i$ el reparto del robot, $\eta=0.46$ la eficiencia nominal y $\ell_i(k)$ la pérdida por perturbación:
$\ell_1=0.35\,d$, $\ell_2=0.175\,d$.

La capacidad de tres piezas no se usa como recorte superior de la planta. Se evalúa como restricción de seguridad mediante
\[
\text{viol}_i(k)=\max\!\left(0,\tilde n_i(k)-3\right),
\]
de forma que un controlador deficiente no pueda ocultar el desbordamiento mediante saturación numérica.

\section{Parámetros del escenario de simulación}

\begin{longtable}{@{}P{0.30\linewidth}P{0.20\linewidth}P{0.38\linewidth}@{}}
\caption{Parámetros del escenario para ambas fases.}
\label{tab:parametros_escenario}\\
\toprule
Parámetro & Valor & Descripción \\ \midrule \endfirsthead
\toprule Parámetro & Valor & Descripción \\ \midrule \endhead
Capacidad $N_{max}$ & 3 piezas & Máximo físico por estera.\\
Nivel objetivo $n^*$ & 1.5 piezas & 50\% capacidad; máxima productividad con margen.\\
Caudal base $\lambda_0$ & 0.42~p/muestra & Producción nominal de la máquina.\\
Pulso de demanda & $+0.35$ en $k=55$--$82$ & Pico de producción.\\
Atasco (cambio peso) & E1: $-35\%$, E2: $-17.5\%$ & Reducción de descarga en $k=95$--$118$.\\
Umbral de riesgo & 2.4 piezas & Desde este nivel el supervisor refuerza el control.\\
Umbral de alarma & 2.7 piezas & Acción de emergencia; reducción de admisión.\\
Límites del actuador & $[0.05,\;1.0]$ (norm.) & Restricción física de velocidad de estera.\\
Condición inicial & $n_1(0)=n_2(0)=1.25$ & Arranque por debajo del objetivo.\\
\bottomrule
\end{longtable}
